\section{범위와 한계}
\label{sec:limitations}

\paragraph{Structural discovery의 범위.}
승격된 structure 결과는 primitive transition에서 observation-only algorithmic
graph discovery를 수행한 뒤 factorized parameter recovery를 적용한다. 따라서
unconstrained neural network가 동일 graph를 스스로 발견한다고 확립하지 않는다.
Neural structure track은 exploratory로 남긴다.

\paragraph{Synthetic bounded-cardinality regime.}
모든 confirmatory world는 선언된 finite synthetic family에서 나온다. Variable-
cardinality audit는 세 개의 finite cardinality panel을 포함하지만 birth, death와
held-out entity identity turnover는 없다. 따라서 zero transition error는 turnover
아래 identity maintenance를 검정하지 않는다.

\paragraph{Temporal validity가 아닌 criterion validity.}
Task-local discrepancy와 regret outcome은 동일한 held-out candidate
trajectory를 사용한다. TSI는 outcome definition에 들어가지 않지만 generic
error comparator와 structural discrepancy를 계산하려면 oracle endpoint가
필요하다. Temporally separated diagnostic의 실패는 현재 benchmark에서 더
강한 prognostic interpretation을 배제한다.

기존 co-primary criterion estimand는 random 및 wrong routing을 포함한 네 arm을
평균한다. Wrong routing이 두 outcome 모두에서 point-effect threshold를 넘으므로
그 analysis는 routing-correctness test가 아니다. Prospective 연구는 direct
correct-versus-shifted routing estimand로 이 특정 ambiguity를 해결하지만,
historical cohort의 estimand나 해석을 바꾸지 않는다.

\paragraph{Comparator dependence.}
Co-primary improvement는 scalar generic model 대비 결과다. Layer-aware
sensitivity는 양수지만 effect floor 미만이다. 따라서 현재 evidence는 common-
correspondence optimization 자체가 모든 단순 layerwise structural
description보다 우월하다고 식별하지 않는다.

\paragraph{Magnitude와 uncertainty.}
Brier improvement는 작다. Simultaneous lower bound는 양수지만 \(0.0025\)
point-effect threshold를 넘지 않는다. Analysis는 positive effect와 threshold
이상의 point estimate를 지지하지만 population effect가 95\% confidence로
threshold를 넘는다고 확립하지 않는다.

\paragraph{External validity와 noise 경계.}
Prospective 결과는 finite state space에서 stochastic coordinate corruption,
세 head family, learned routing과 held-out two-mechanism composition을 다룬다.
Project import가 없는 Node.js code path가 120개 learned factorized fit과 보고한
analysis mean 열 개를 모두 재현한다. 이는 independent implementation check이지
independent research-group study가 아니다. Visual input, entity turnover와
real-world system은 정의한 population 밖이며 이 결과에서 추론하지 않는다.
